\chapter{Audit of the frozen puzzle and modulus models}
\label{chap:frozen-audit}

The purpose of a formal proof is to check a mathematical statement with
its intended definitions.  Checking an implication and identifying it
with a named theorem are separate obligations.  This chapter makes the
distinction concrete for the inherited puzzle layer.  The obstruction
below is proved in Lean; the dependency itself is left pinned and
unchanged so that the audit remains reproducible.

\section{Two different quantities called modulus}

The classical conformal modulus of the round annulus
\[
 A(r,R)=\{z:r<|z|<R\},\qquad 0<r<R,
\]
is
\[
 \modulus A(r,R)=\frac{1}{2\pi}\log\frac Rr.
\]
It is invariant under conformal isomorphisms, in particular under a
simultaneous rescaling of $r$ and $R$.

In the pinned Yoccoz dependency, however, the declarations
\lean{MLC.Quadratic.weight} and \lean{MLC.Quadratic.modulus} are
\[
 w(z)=e^{-|z|^2},\qquad
 \weightedarea(A)=\int_A e^{-|z|^2}\,dA(z).
\]
The defining file is
\lean{Yoccoz/Quadratic/Complex/Groetzsch.lean} at the revision recorded
in \lean{lake-manifest.json}.  Its own comment describes the quantity
as a proxy.  The formalization must preserve that qualification.

Weighted area is additive on disjoint measurable sets and is bounded by
the finite integral over $\C$.  It is not conformally invariant: moving
a fixed disk far from the origin makes its weighted area tend to zero.
Consequently, a finite-union identity for this quantity does not prove a
classical conformal-modulus identity.

\begin{proposition}[A round-annulus comparison]
For $A_n=A(2^{-(n+1)},2^{-n})$,
\[
 \modulus(A_n)=\frac{\log2}{2\pi},
\]
whereas
\[
 \weightedarea(A_n)
 =\pi\left(e^{-4^{-(n+1)}}-e^{-4^{-n}}\right).
\]
Thus the conformal-modulus series diverges and the weighted-area series
converges.
\end{proposition}

\begin{proof}
The conformal formula follows from the ratio of the radii.
Polar coordinates give
\[
 \weightedarea(A(r,R))
 =2\pi\int_r^R e^{-t^2}t\,dt
 =\pi(e^{-r^2}-e^{-R^2}).
\]
The displayed weighted-area series telescopes to
$\pi(1-e^{-1})$.  Its convergence is compatible with the singleton
intersection of the nested closed disks.
\end{proof}

This example shows that even the simplest classical shrinking mechanism
cannot be recovered by replacing conformal modulus with the present
weighted area.

\section{A universal summability obstruction}

\begin{theorem}[Weighted areas of nested differences]
\label{thm:weighted-summability}
Let $(P_n)$ be any nested sequence of measurable subsets of $\C$.
Then
\[
 \sum_{n\ge0}\weightedarea(P_n\setminus P_{n+1})<\infty.
\]
Neither connectedness nor a hypothesis on the intersection is needed.
\end{theorem}

\begin{proof}
The sets $A_n=P_n\setminus P_{n+1}$ are pairwise disjoint.
Indeed, if $i<j$, then $P_j\subseteq P_{i+1}$, so a point of $A_j$
cannot lie in $A_i$.  Finite additivity and nonnegativity give
\[
 0\le\sum_{n<N}\weightedarea(A_n)
   =\weightedarea\left(\bigcup_{n<N}A_n\right)
   \le\weightedarea(\C)<\infty.
\]
The nonnegative partial sums are bounded, so the series converges.
\end{proof}

\paragraph{Lean correspondence.}
The bound and summability theorem are
\leanref{Mlc/Literature/ModelAudit.lean}{weightedArea_le_total} and
\leanref{Mlc/Literature/ModelAudit.lean}{summable_weightedArea_differences}.
The proof uses the dependency's actual integral definition and its
integrability theorem.  It introduces no conformal-modulus axiom.

The frozen pieces are
\[
 P_n^{\rm fr}(c)=
 \operatorname{Comp}_0\{z:G_c(z)<2^{-n}\}.
\]
They are components of full Green sublevels, with no landing-ray graph
removed.  They are open and nested.  Hence
Theorem~\ref{thm:weighted-summability} applies to their differences for
every parameter, including parameters outside $\M$.

\begin{corollary}[The legacy finite class is empty]
\label{cor:empty-frozen-class}
For every $c\in\C$, the series
\[
 \sum_n\weightedarea
       (P_n^{\rm fr}(c)\setminus P_{n+1}^{\rm fr}(c))
\]
is summable.  Consequently the predicate currently named
\lean{MLC.FinitelyRenormalizable} is false for every $c$.
\end{corollary}

\begin{proof}
The first assertion is the preceding theorem.  The second follows by
unfolding the definitions: \lean{FinitelyRenormalizable} abbreviates
\lean{NonRenormalizable}, and the latter is precisely the negation
of summability of this weighted-area sequence.
\end{proof}

\begin{samepage}
The exact declarations are
\begin{center}\small
\leanref{Mlc/Literature/ModelAudit.lean}{frozen_puzzle_modulus_summable}\\
\leanref{Mlc/Literature/ModelAudit.lean}{not_finitelyRenormalizable_in_frozenModel}.
\end{center}
\end{samepage}
Thus the implication named \lean{yoccoz_puzzle_shrinkage} is valid
as a Lean proposition, but its premise has no witnesses.  The same is
true of family statements indexed by this legacy predicate.  A
successful axiom audit alone would not reveal this semantic problem.

\section{The principal-nest certificate is also obstructed}

Taking a monotone subsequence $P_{d_n}^{\rm fr}(c)$ gives another
nested measurable family.  Its successive differences therefore have
summable weighted areas, independently of how the depths were selected.

\begin{corollary}[No positive uniform proxy bound]
\label{cor:no-frozen-certificate}
There is no $\mu>0$ uniformly bounded above by all the weighted areas
of the frozen puzzle annuli.  The analogous lower-bound field of
\lean{SatelliteAPrioriCertificate} also has no instances.
\end{corollary}

\begin{proof}
A summable real sequence tends to zero, which excludes a fixed positive
lower bound.  Apply this first to all consecutive depths and then to the
monotone depths of the proposed satellite tower.
\end{proof}

The formal statements are
\lean{no_uniform_lower_bound_frozen_puzzleAnnuli},
\lean{frozen_principalNest_modulus_summable}, and
\leanref{Mlc/Literature/ModelAudit.lean}{not_nonempty_satelliteAPrioriCertificate}.
This is a refuted input with the present definitions.  It cannot be
listed merely as an estimate that might later be supplied by a
published a priori-bound theorem.

\section{Other names that require a semantic audit}

The legacy declarations retain their names for compatibility with the
current root-facing code.  Their mathematical meanings are:

\begin{center}\small
\begin{tabularx}{\textwidth}{@{}P{0.38\textwidth}Y@{}}
\textbf{Declaration} & \textbf{Actual content}\\
\hline
\lean{PrimitiveRenormalizable} &
The pointwise connected-neighbourhood conclusion itself, quantified
over membership in $\M$; no primitive polynomial-like restriction is
constructed.\\
\lean{FinitelyRenormalizable} &
Nonsummability of frozen Gaussian annulus areas; the predicate is empty.\\
\lean{InfinitelyRenormalizable} &
Summability of the same sequence, so its field is satisfiable for every
parameter; it does not classify genuine infinite renormalization.\\
\lean{ParaPuzzlePieceAt} &
A translation of a dynamical Green-sublevel component, rather than a
parameter cell defined by moving landing rays.\\
\lean{cumulativeDepth} &
An additive indexing sum; original return periods instead multiply.\\
\end{tabularx}
\end{center}

These observations concern the formal models.  They do not refute
Yoccoz's, Lyubich's, or Dudko--Lyubich's results.
The generic \lean{ParameterPuzzleBasis} and
\lean{LiteratureMLCData} interfaces avoid prescribing the frozen sets,
so their topological endpoint theorems remain useful.  They still need
constructors from the actual source hypotheses.

\section{What remains sound, and what must be replaced}

The escape-radius identities, compact outer approximation, component
intersection theorem, uniform-buffer implication, and continuous
retraction theorem use their stated orbit and topological definitions.
Their proofs do not turn Gaussian area into conformal modulus.
Likewise, categorical image/intersection identities express genuine
topological facts.  These remain available as endpoints for future
analytic constructions.

\begin{formalobligation}[A genuine conformal puzzle layer]
Introduce ring domains with their boundary markings and conformal modulus;
prove its invariance, the serial Grötzsch inequality, and the shrinking
criterion.  Construct landing-ray graph cuts and their pullback puzzle
domains.  Define finite and infinite renormalization by marked
polynomial-like restrictions, and prove the geometric estimates from
those definitions.  Finally realize the associated parameter pieces and
apply Theorem~\ref{thm:parameter-shrinking}.
\end{formalobligation}

The analytic chapter states the classical objects and results with these
meanings.  The conditional theorem matrix records which current Lean
arrows are model-independent and which have refuted frozen premises.
This keeps a readable mathematical narrative while exposing the exact
distance between the existing code and the published literature.
